\clearpage\gdef\otChapNum{0}\pagestyle{otFront}\phantomsection\label{couverture-programme}\pdfbookmark[0]{Couverture du programme}{bm-programme}
\otPageTitre{Repères du programme}{Contenus, méthodes et exercices correspondants}
Le découpage ci-dessous est celui du livre. Les trois ressources transversales se réinvestissent dans tous les thèmes ; leur regroupement n’impose pas d’attendre la fin de l’année. Les références renvoient à un cours et à un exemple d’exercice, sans limiter les notions à ces seules pages.\par\medskip
\begingroup\fontsize{9.5}{12}\selectfont\renewcommand{\arraystretch}{1.2}
\begin{longtable}{@{}p{106mm}p{53mm}@{}}\toprule
Contenu ou capacité travaillé & Où le retrouver\\\midrule\endhead
\multicolumn{2}{@{}l}{\otSemi\color{otNavy}1. Arithmétique et nombres entiers}\\[2mm]
Multiples, diviseurs et parité & \hyperref[ess-1]{Cours p.\,\pageref*{ess-1}} ; \hyperref[exo-1-1]{ex. 1}\\[2mm]
Fractions irréductibles & \hyperref[ess-1]{Cours p.\,\pageref*{ess-1}} ; \hyperref[exo-1-4]{ex. 4}\\[2mm]
Somme de multiples et carré impair & \hyperref[ess-1]{Cours p.\,\pageref*{ess-1}} ; \hyperref[exo-1-6]{ex. 6}\\[2mm]
Algorithmes de divisibilité & \hyperref[ess-1]{Cours p.\,\pageref*{ess-1}} ; \hyperref[exo-1-8]{ex. 8}\\[2mm]
\multicolumn{2}{@{}l}{\otSemi\color{otNavy}2. Réels et intervalles}\\[2mm]
Ensembles de nombres & \hyperref[ess-2]{Cours p.\,\pageref*{ess-2}} ; \hyperref[exo-2-1]{ex. 1}\\[2mm]
Intervalles, union et intersection & \hyperref[ess-2]{Cours p.\,\pageref*{ess-2}} ; \hyperref[exo-2-4]{ex. 4}\\[2mm]
Valeur absolue et distance & \hyperref[ess-2]{Cours p.\,\pageref*{ess-2}} ; \hyperref[exo-2-6]{ex. 6}\\[2mm]
Encadrement et chiffres significatifs & \hyperref[ess-2]{Cours p.\,\pageref*{ess-2}} ; \hyperref[exo-2-10]{ex. 10}\\[2mm]
\multicolumn{2}{@{}l}{\otSemi\color{otNavy}3. Calcul numérique et littéral}\\[2mm]
Puissances et racines & \hyperref[ess-3]{Cours p.\,\pageref*{ess-3}} ; \hyperref[exo-3-2]{ex. 2}\\[2mm]
Identités remarquables & \hyperref[ess-3]{Cours p.\,\pageref*{ess-3}} ; \hyperref[exo-3-5]{ex. 5}\\[2mm]
Expressions fractionnaires et domaines & \hyperref[ess-3]{Cours p.\,\pageref*{ess-3}} ; \hyperref[exo-3-7]{ex. 7}\\[2mm]
Isolement de variables et comparaison & \hyperref[ess-3]{Cours p.\,\pageref*{ess-3}} ; \hyperref[exo-3-9]{ex. 9}\\[2mm]
\multicolumn{2}{@{}l}{\otSemi\color{otNavy}4. Équations et inéquations}\\[2mm]
Équations linéaires, produit et quotient & \hyperref[ess-4]{Cours p.\,\pageref*{ess-4}} ; \hyperref[exo-4-6]{ex. 6}\\[2mm]
Inéquations et règles de l’ordre & \hyperref[ess-4]{Cours p.\,\pageref*{ess-4}} ; \hyperref[exo-4-3]{ex. 3}\\[2mm]
Tableaux de signes & \hyperref[ess-4]{Cours p.\,\pageref*{ess-4}} ; \hyperref[exo-4-8]{ex. 8}\\[2mm]
Modélisation et paramètres & \hyperref[ess-4]{Cours p.\,\pageref*{ess-4}} ; \hyperref[exo-4-14]{ex. 14}\\[2mm]
\multicolumn{2}{@{}l}{\otSemi\color{otNavy}5. Fonctions et représentations}\\[2mm]
Images, antécédents et courbes & \hyperref[ess-5]{Cours p.\,\pageref*{ess-5}} ; \hyperref[exo-5-4]{ex. 4}\\[2mm]
Ensembles de définition & \hyperref[ess-5]{Cours p.\,\pageref*{ess-5}} ; \hyperref[exo-5-6]{ex. 6}\\[2mm]
Équations et inéquations graphiques & \hyperref[ess-5]{Cours p.\,\pageref*{ess-5}} ; \hyperref[exo-5-9]{ex. 9}\\[2mm]
Modèles continus et discrets & \hyperref[ess-5]{Cours p.\,\pageref*{ess-5}} ; \hyperref[exo-5-12]{ex. 12}\\[2mm]
\multicolumn{2}{@{}l}{\otSemi\color{otNavy}6. Les fonctions de référence}\\[2mm]
Fonctions affine, valeur absolue, carré, cube, inverse et racine & \hyperref[ess-6]{Cours p.\,\pageref*{ess-6}} ; \hyperref[exo-6-1]{ex. 1}\\[2mm]
Domaines, signes et variations & \hyperref[ess-6]{Cours p.\,\pageref*{ess-6}} ; \hyperref[exo-6-8]{ex. 8}\\[2mm]
Équations et inéquations de référence & \hyperref[ess-6]{Cours p.\,\pageref*{ess-6}} ; \hyperref[exo-6-7]{ex. 7}\\[2mm]
Position de x et de son carré & \hyperref[ess-6]{Cours p.\,\pageref*{ess-6}} ; \hyperref[exo-6-9]{ex. 9}\\[2mm]
\multicolumn{2}{@{}l}{\otSemi\color{otNavy}7. Variations et optimisation}\\[2mm]
Monotonie et extrémums & \hyperref[ess-7]{Cours p.\,\pageref*{ess-7}} ; \hyperref[exo-7-6]{ex. 6}\\[2mm]
Variations des affines & \hyperref[ess-7]{Cours p.\,\pageref*{ess-7}} ; \hyperref[exo-7-7]{ex. 7}\\[2mm]
Optimisation sans dérivée & \hyperref[ess-7]{Cours p.\,\pageref*{ess-7}} ; \hyperref[exo-7-11]{ex. 11}\\[2mm]
Balayage et preuve algébrique & \hyperref[ess-7]{Cours p.\,\pageref*{ess-7}} ; \hyperref[exo-7-10]{ex. 10}\\[2mm]
\multicolumn{2}{@{}l}{\otSemi\color{otNavy}8. Vecteurs du plan}\\[2mm]
Coordonnées, norme et opérations vectorielles & \hyperref[ess-8]{Cours p.\,\pageref*{ess-8}} ; \hyperref[exo-8-5]{ex. 5}\\[2mm]
Relation de Chasles et milieu & \hyperref[ess-8]{Cours p.\,\pageref*{ess-8}} ; \hyperref[exo-8-14]{ex. 14}\\[2mm]
Déterminant et colinéarité & \hyperref[ess-8]{Cours p.\,\pageref*{ess-8}} ; \hyperref[exo-8-7]{ex. 7}\\[2mm]
Décomposition dans une base & \hyperref[ess-8]{Cours p.\,\pageref*{ess-8}} ; \hyperref[exo-8-13]{ex. 13}\\[2mm]
\multicolumn{2}{@{}l}{\otSemi\color{otNavy}9. Repérage et configurations}\\[2mm]
Distances et coordonnées du milieu & \hyperref[ess-9]{Cours p.\,\pageref*{ess-9}} ; \hyperref[exo-9-2]{ex. 2}\\[2mm]
Configurations planes et médiatrices & \hyperref[ess-9]{Cours p.\,\pageref*{ess-9}} ; \hyperref[exo-9-5]{ex. 5}\\[2mm]
Réactivation de Pythagore, Thalès et trigonométrie & \hyperref[ess-9]{Cours p.\,\pageref*{ess-9}} ; \hyperref[exo-9-8]{ex. 8}\\[2mm]
Aires, volumes et réductions & \hyperref[ess-9]{Cours p.\,\pageref*{ess-9}} ; \hyperref[exo-9-9]{ex. 9}\\[2mm]
\multicolumn{2}{@{}l}{\otSemi\color{otNavy}10. Droites du plan}\\[2mm]
Équations réduite et cartésienne & \hyperref[ess-10]{Cours p.\,\pageref*{ess-10}} ; \hyperref[exo-10-7]{ex. 7}\\[2mm]
Vecteur directeur et pente & \hyperref[ess-10]{Cours p.\,\pageref*{ess-10}} ; \hyperref[exo-10-6]{ex. 6}\\[2mm]
Alignement et parallélisme & \hyperref[ess-10]{Cours p.\,\pageref*{ess-10}} ; \hyperref[exo-10-13]{ex. 13}\\[2mm]
Systèmes et intersections & \hyperref[ess-10]{Cours p.\,\pageref*{ess-10}} ; \hyperref[exo-10-9]{ex. 9}\\[2mm]
\multicolumn{2}{@{}l}{\otSemi\color{otNavy}11. Proportions et évolutions}\\[2mm]
Proportion et pourcentages emboîtés & \hyperref[ess-11]{Cours p.\,\pageref*{ess-11}} ; \hyperref[exo-11-6]{ex. 6}\\[2mm]
Variations additive et multiplicative & \hyperref[ess-11]{Cours p.\,\pageref*{ess-11}} ; \hyperref[exo-11-4]{ex. 4}\\[2mm]
Évolutions successives et réciproques & \hyperref[ess-11]{Cours p.\,\pageref*{ess-11}} ; \hyperref[exo-11-7]{ex. 7}\\[2mm]
Lecture critique de données réelles & \hyperref[ess-11]{Cours p.\,\pageref*{ess-11}} ; \hyperref[exo-11-14]{ex. 14}\\[2mm]
\multicolumn{2}{@{}l}{\otSemi\color{otNavy}12. Statistique descriptive}\\[2mm]
Moyenne, médiane et quartiles & \hyperref[ess-12]{Cours p.\,\pageref*{ess-12}} ; \hyperref[exo-12-2]{ex. 2}\\[2mm]
Écart type et transformations de données & \hyperref[ess-12]{Cours p.\,\pageref*{ess-12}} ; \hyperref[exo-12-7]{ex. 7}\\[2mm]
Histogrammes et polygones cumulés & \hyperref[ess-12]{Cours p.\,\pageref*{ess-12}} ; \hyperref[exo-12-10]{ex. 10}\\[2mm]
Moyenne et médiane pour des classes & \hyperref[ess-12]{Cours p.\,\pageref*{ess-12}} ; \hyperref[exo-12-14]{ex. 14}\\[2mm]
\multicolumn{2}{@{}l}{\otSemi\color{otNavy}13. Tableaux croisés et fréquences}\\[2mm]
Variables nominales et ordinales & \hyperref[ess-13]{Cours p.\,\pageref*{ess-13}} ; \hyperref[exo-13-5]{ex. 5}\\[2mm]
Tableaux croisés et marges & \hyperref[ess-13]{Cours p.\,\pageref*{ess-13}} ; \hyperref[exo-13-1]{ex. 1}\\[2mm]
Fréquences conditionnelles & \hyperref[ess-13]{Cours p.\,\pageref*{ess-13}} ; \hyperref[exo-13-2]{ex. 2}\\[2mm]
Filtres et lecture de données réelles & \hyperref[ess-13]{Cours p.\,\pageref*{ess-13}} ; \hyperref[exo-13-10]{ex. 10}\\[2mm]
\multicolumn{2}{@{}l}{\otSemi\color{otNavy}14. Probabilités et arbres}\\[2mm]
Univers, équiprobabilité et événements & \hyperref[ess-14]{Cours p.\,\pageref*{ess-14}} ; \hyperref[exo-14-1]{ex. 1}\\[2mm]
Conditionnement et inversion & \hyperref[ess-14]{Cours p.\,\pageref*{ess-14}} ; \hyperref[exo-14-11]{ex. 11}\\[2mm]
Arbres et probabilités de chemin & \hyperref[ess-14]{Cours p.\,\pageref*{ess-14}} ; \hyperref[exo-14-7]{ex. 7}\\[2mm]
Simulation et fluctuation & \hyperref[ess-14]{Cours p.\,\pageref*{ess-14}} ; \hyperref[exo-14-14]{ex. 14}\\[2mm]
\multicolumn{2}{@{}l}{\otSemi\color{otNavy}15. Algorithmique et programmation}\\[2mm]
Types, affectations, conditions et boucles & \hyperref[ess-15]{Cours p.\,\pageref*{ess-15}} ; \hyperref[exo-15-2]{ex. 2}\\[2mm]
Fonctions et tests & \hyperref[ess-15]{Cours p.\,\pageref*{ess-15}} ; \hyperref[exo-15-3]{ex. 3}\\[2mm]
Simulation, moyenne et écart type & \hyperref[ess-15]{Cours p.\,\pageref*{ess-15}} ; \hyperref[exo-15-12]{ex. 12}\\[2mm]
Algorithmes géométriques, balayage et dichotomie & \hyperref[ess-15]{Cours p.\,\pageref*{ess-15}} ; \hyperref[exo-15-15]{ex. 15}\\[2mm]
\multicolumn{2}{@{}l}{\otSemi\color{otNavy}16. Raisonner et faire la synthèse}\\[2mm]
Ensembles, couples et cardinal & \hyperref[ess-16]{Cours p.\,\pageref*{ess-16}} ; \hyperref[exo-16-2]{ex. 2}\\[2mm]
Connecteurs, implication et réciproque & \hyperref[ess-16]{Cours p.\,\pageref*{ess-16}} ; \hyperref[exo-16-4]{ex. 4}\\[2mm]
Disjonction, contraposée et absurde & \hyperref[ess-16]{Cours p.\,\pageref*{ess-16}} ; \hyperref[exo-16-8]{ex. 8}\\[2mm]
Problèmes de synthèse & \hyperref[ess-16]{Cours p.\,\pageref*{ess-16}} ; \hyperref[exo-16-11]{ex. 11}\\[2mm]\bottomrule\end{longtable}\endgroup
\par\medskip\otIntertitre{Les trois parties transversales}
{\fontsize{9.5}{12}\selectfont\renewcommand{\arraystretch}{1.2}\begin{tabularx}{\linewidth}{@{}Yl@{}}\toprule
Vocabulaire ensembliste et logique : opérations sur les ensembles, cardinal, propositions et raisonnements & \hyperref[ess-16]{Cours p.\,\pageref*{ess-16}} ; \hyperref[exo-16-1]{ex. 1}\\[1mm]
Algorithmique : affectation notée $\leftarrow$, langage naturel et Python, fonctions, hasard & \hyperref[ess-15]{Cours p.\,\pageref*{ess-15}} ; \hyperref[exo-15-1]{ex. 1}\\[1mm]
Automatismes : 30 capacités du programme, 40 questions flash corrigées & \hyperref[automatismes-1]{Séries p.\,\pageref*{automatismes-1}}\\\bottomrule
\end{tabularx}}\par\smallskip
{\fontsize{9.5}{12}\selectfont Nuage de points : atelier de lecture graphique, p.\,\pageref{atelier-nuage}. Signe d’une fonction affine et taux d’accroissement : cours du chapitre~6.}
\clearpage\otPageTitre{Preuves et prolongements}{Un périmètre explicite}
\otIntertitre{Les démonstrations identifiées dans le programme}
Somme de multiples et carré impair ; non-décimalité de $1/3$ et irrationalité de $\sqrt2$ ; racine d’un produit ; critère de colinéarité ; équation cartésienne de droite ; variations des affines ; position de $x$ et $x^2$ ; variations du carré et de l’inverse. Elles sont rassemblées dans les 13 démonstrations de référence, avec deux compléments.\par\medskip
\otIntertitre{Toutes les pistes d’approfondissement proposées}
\begin{tabularx}{\linewidth}{@{}Yl@{}}\toprule
Écritures décimales illimitées ; périodicité rationnelle & D01\\[1.5mm]
Carré de trois termes ; cube d’une somme & D02\\[1.5mm]
Moyennes arithmétique et géométrique & D03\\[1.5mm]
Barycentres de deux et trois points & D04\\[1.5mm]
Formule des coordonnées du milieu, démontrée & Preuve 12\\[1.5mm]
Concours des hauteurs & D06\\[1.5mm]
Aire $ab\sin C/2$ & D07\\[1.5mm]
Isobarycentre et concours des médianes & D05\\[1.5mm]
Médiatrices et cercle circonscrit & D08\\[1.5mm]
Équidistance d’un point et de l’axe horizontal & D09\\[1.5mm]
Régions décrites par des inégalités & D10\\[1.5mm]
Lien entre courbes carré et racine sur les positifs & D11\\\bottomrule
\end{tabularx}\par\medskip
Les défis D12 à D16 ajoutent d’autres recherches : paramètre, dispersion, marges d’un tableau, contrôle et dénombrement. Le caractère facultatif est indiqué dans le parcours.\par\medskip
\otIntertitre{Les exemples d’algorithmes}
Divisibilité et dernier multiple : 1.8 et 1.9. Encadrement de racine : 2.15 et 15.7. Première puissance au-delà ou en deçà d’un seuil : 11.13 et 15.6. Alignement et équation de droite : 15.8 et 15.9. Balayage et dichotomie d’un extrémum : 7.10, 15.15 et atelier. Longueur approchée de courbe : 15.14. Filtres ET, OU, NON sur deux listes, tableau croisé et fréquences conditionnelles : 13.7, 13.8 et atelier. Simulation et loi des grands nombres : 15.10 à 15.12.
\clearpage